\section{Conclusion and Future Scope}

\subsection{Conclusion}
The college timetable is consulted by everyone every day, yet it exists only as a dense PDF that has
to be read by hand, with no search, no personal view, no version control and no clash checking. This
mini project built \emph{Dynamic Timetable}, a role-based timetable monitoring system that reads the
timetable as it is actually published (PDF, scanned image, Excel or Word), extracts every class with
its source page and a verification status, detects clashes, and serves students, professors, the HOD,
the Principal and the Admin through one app on Android, iPhone and laptop. Users can ask everyday
questions in plain English. A deterministic compositional planner turns each question into a typed
query plan (operation, filters, lab/lecture, date or week, time window, ordering) and compiles it into
parameterised queries, so it answers counts, first/last and day-wise questions as well as the everyday
ones in about 40~ms, never generates SQL from text, and asks for clarification instead of guessing. On the college's real 10-page timetable the system verified 452 of 453 class entries
automatically and found the two genuine clashes in the published document. The system is covered by
168 back-end tests, 20 app tests and end-to-end browser tests, and it runs as a free online demo.

\subsection{Achievement of Objectives}
\begin{itemize}
\item \textbf{Achieved:} O1 (99.8\% verified), O2 (all formats through one pipeline), O3 (five roles,
      server-side), O4 (composed query plans plus specialised operations), O5 (0 false clash reports), O6 (strict domain separation and
      archives), O9 (all tests passing).
\item \textbf{Achieved on development hardware:} O7 (latency and correctness under 32 concurrent
      clients). It still has to be confirmed on production hardware.
\item \textbf{Partially achieved:} O8. The app runs on Android and iPhone through Expo Go and on
      laptops through the web and the online demo, but installable store builds (APK/TestFlight) have
      not yet been produced.
\end{itemize}

\subsection{Major Engineering Contributions}
\begin{enumerate}
\item A timetable-aware extraction pipeline with constrained backtracking for inconsistent time labels,
      stacked lab-group parsing, page-local legends and per-entry provenance.
\item A conflict detector that separates intended overlaps (parallel labs, electives, shared sessions)
      from genuine clashes in $O(n\log n + m)$.
\item A safe, compositional natural-language interface: questions become typed query plans (operation,
      filters, lab/lecture, date or week, time window, ordering) compiled into parameterised queries with
      timetable-scoped entity resolution and clarification instead of guessing, injection-proof by design
      and no longer limited to a fixed list of questions.
\item A secure multi-role architecture: Argon2id, rotating refresh tokens with reuse detection,
      database-level ownership constraints (composite FKs, partial unique indexes) and audited corrections.
\item A free, reproducible deployment: Docker image, infrastructure as code (\texttt{render.yaml}) and
      continuous deployment from Git.
\end{enumerate}

\subsection{Current Limitations}
\begin{itemize}
\item \textbf{Technical:} extraction accuracy has not been measured against a hand-labelled ground truth;
      new timetable layouts may need parser rules; the rate limiter is per process; co-taught cells
      (``RRS+SB'') are not linked to both teachers' personal schedules.
\item \textbf{Functional:} the system does not generate timetables; questions outside the 30
      intents are politely refused; teacher location is the \emph{scheduled} room, not tracked.
\item \textbf{Operational:} the free demo sleeps when idle and its database expires every 30 days;
      push notifications and e-mail are not yet active; room floors were derived from room numbers and
      must be checked by the Admin.
\end{itemize}

\subsection{Future Enhancements}
\begin{enumerate}
\item Measure precision/recall on hand-labelled pages and publish per-field accuracy.
\item College single sign-on (OIDC/Google Workspace) instead of separate passwords.
\item Store builds (Android APK / Play Store, iOS TestFlight) with push notifications for timetable
      changes and ``class starting in 10 minutes'' reminders.
\item ``What changed?'' diff between two official versions, sent to affected users only.
\item Bulk CSV import of students, staff and rooms; many-to-many teacher links for co-taught classes.
\item An optional LLM \emph{suggester} that proposes an intent and parameters for unusual phrasings,
      which must still pass the same validation and templates (never free SQL).
\item Room-utilisation analytics for the HOD and Principal; exam-timetable support.
\item Capacity testing on production hardware; a shared rate-limit store and vocabulary caching for
      multi-instance deployment.
\end{enumerate}

\subsection{Lessons Learned}
\begin{itemize}
\item \textbf{Test with real data early.} The real PDF broke assumptions that synthetic data never would.
\item \textbf{Make uncertainty visible.} Statuses, raw text and page numbers turned possible silent
      errors into review tasks.
\item \textbf{Security is a design decision.} Choosing templates over generated SQL removed a whole
      class of attacks; our own audit still found a critical bug (P0-1), which showed us the value of reviewing our own code adversarially.
\item \textbf{Automated tests make change safe.} The SDK upgrade and the parser changes were made with
      confidence because 200+ tests guarded behaviour.
\item \textbf{Project management:} risk-first ordering (extraction before UI), small increments,
      one commit per working increment, and documenting decisions (why, not only what) kept the
      project on schedule and easy to explain.
\end{itemize}

\newpage
